\section{D. Splittings, their ambiguity, and the sign of the norm}

\begin{claim}{D.1}{Known}{The extension splits, with a unique Belyi normalization}
The \S46 extension in question is
\[
 1\longrightarrow\widehat F\longrightarrow\widetilde M^!
 \longrightarrow\widehat{\mathrm{GT}}_0\longrightarrow1,
\]
not an extension denoted $\widetilde N$ by $\widehat F$. It has the section whose representative is
$x\mapsto x^\mu$, $y\mapsto f^{-1}y^\mu f$, $f\in\widehat F'$. That section is unique among sections all of whose representatives have this normalization. It is Harbater--Schneps (2000), \S1.2, Proposition 2 and Theorem 1, applied through A.3. The statements of uniqueness on scans208--209 concern the further matrix normalizer or torus constraint and are accompanied by explicit withdrawal margins.
\end{claim}
% anchor:D.1
\manuscriptanchor{\S 46; fols403--403 verso; scans208--209}

\begin{proof}
We give the normalization argument, since it identifies the ambiguity. Start with a loop automorphism and conjugate it so that it sends $x$ to $x^\mu$. Write its image of $y$ as $g^{-1}y^\mu g$. Multiplying $g$ on the left by a power of $y$ does not change this image, and composing the automorphism with an inner power of $x$ changes its right $x$ exponent. These two changes kill both abelianization coordinates of $g$. If two normalized representatives differ by $\operatorname{Int}(h)$, comparison on $x$ gives $h=x^c$; comparison on $y$ then gives $g x^c f^{-1}=y^d$. Abelianization and $f,g\in\widehat F'$ give $c=d=0$, so both representatives and both words coincide. The composite of two normalized representatives is again normalized: its conjugator is a product of commutator-subgroup elements, since that subgroup is characteristic. Thus the unique choice is a group homomorphism, and the closed conditions defining its graph make it continuous. This is also the proof of the cited theorem in the present coordinates.

The scan209 ink names $\widetilde M^!$ and says ``Je suspecte que c'est le seul scindage''; the margin reads ``canul\'e, pour m\^emes raisons que tout \`a l'heure''. The preceding scan208 already retracts its analogous claim with an isogeny objection. We therefore distinguish these withdrawn conjectures from the theorem of unique normalized representatives.
\end{proof}

\begin{claim}{D.2}{Refuted}{Unqualified strict uniqueness of a splitting}
There is more than one strict continuous section of the displayed extension. The strongest replacement furnished by the equations is uniqueness after the Belyi normalization of D.1, or a classification by a nonabelian continuous $H^1$ after quotienting by conjugation. In particular, ``unique'' without naming the normalization or the conjugation ambiguity is false.
\end{claim}
% anchor:D.2
\manuscriptanchor{\S 46; fols403--403 verso; scans208--209}

\begin{proof}
Let $s$ be the normalized section and $c=(-1,1)$. For any $k\in\widehat F$ put $s_k(g)=k s(g)k^{-1}$. This is a section because $k$ lies in the geometric kernel. For $k=x$, since $s(c)(x)=x^{-1}$,
\[
 s_x(c)=\operatorname{Int}(x^2)\,s(c)\ne s(c).
\]
The inequality follows from $x^2\ne1$ and the trivial center of $\widehat F$. Thus the sections differ; they are nevertheless conjugate by construction. For a general section $t$, the quotient $a(g)=t(g)s(g)^{-1}$ belongs to $\widehat F$ and satisfies
\[
 a(gh)=a(g)\,s(g)(a(h)).
\]
Conversely this cocycle equation makes $a(g)s(g)$ a section. Conjugating by $k$ changes $a(g)$ to $k a(g)s(g)(k^{-1})$. These are exactly the defining equivalence relations of $H^1_{\mathrm{cts}}(\widehat{\mathrm{GT}}_0,\widehat F)$ for the action determined by $s$.
\end{proof}

\begin{claim}{D.3}{Refuted}{The strict normalizer uniqueness already fails modulo seven}
The matrix normalizer condition on scan208 does not restore strict uniqueness. Let $K_{1/2}$ be the section centralizing $\sigma_\infty$. There is a distinct conjugate section whose image still maps into the normalizer of $\widehat{\mathbb Z}[S]^\times$ in the projective matrix group. Distinctness is detected by the two involutions
\[
 J=\begin{pmatrix}1&0\\0&6\end{pmatrix},\qquad
 J'=\begin{pmatrix}0&1\\1&0\end{pmatrix}
 \quad\text{in }\operatorname{PGL}_2(\mathbb F_7).
\]
This is a counterexample to strict uniqueness; it does not refute uniqueness modulo geometric conjugation or the distinct order-three torus question of scan209.
\end{claim}
% anchor:D.3
\manuscriptanchor{\S 46; fol403; scan208}

\begin{proof}
The cyclic finite-subgroup theorem for free profinite products says that an order-two element of $\widehat B=\widehat C_2\amalg\widehat C_3$ is conjugate to $s$, with centralizer $\langle s\rangle$. Applied to the image of $s$ under a pure automorphism, its conjugator may be taken in $\widehat F$: its reduction in $S_3$ centralizes the transposition and is removed by a power of $s$. The intersection of this centralizer with $\widehat F$ is trivial. Consequently each class has a unique representative centralizing $s$, so these representatives form the continuous section $K_{1/2}$. This is the precise profinite input behind the scan206--208 construction. The torsion part is Herfort--Ribes (1985), Theorem 2(i), explicitly cited in the opened Minamide--Nakamura paper, p.8; the primary paper was not supplied and its proof has not been independently checked here. The normal-form centralizer assertion is the corresponding free-product centralizer theorem. We isolate this standard structural input instead of inferring it from a finite matrix test.

Here is the counterexample once that section is fixed. At $7$ set
\[
 k_7=2I+2S=\begin{pmatrix}2&5\\2&2\end{pmatrix}.
\]
It has determinant one, lies in the $S$-torus, and $k_7Jk_7^{-1}=J'$. Both $J,J'$ square to $I$ and conjugate $S$ to $-S$. For a global lift, take $t_7=2$ and the unique Hensel lift $u_7\in\mathbb Z_7$ of $2$ with $u_7^2=-3$; its derivative $2u_7$ is a unit, so the lift exists. Set $k_7^{\mathrm{ad}}=u_7I+t_7S$ and $k_p^{\mathrm{ad}}=I$ for all $p\ne7$. Then $k^{\mathrm{ad}}\in\operatorname{SL}_2(\widehat{\mathbb Z})$, is the identity modulo two, and commutes with $S$.

Reduction $\widehat F\to\ker((\operatorname{SL}_2(\widehat{\mathbb Z})/\{\pm I\})\to\operatorname{SL}_2(\mathbb F_2))$ is onto. Indeed $\operatorname{SL}_2(\mathbb Z)$ surjects onto $\operatorname{SL}_2(\mathbb Z/n)$ by elementary row operations and the Chinese remainder theorem; restricting to the inverse image of the identity mod two and then taking compact inverse limits proves the assertion. Hence choose $k\in\widehat F$ over $k^{\mathrm{ad}}$. Conjugate $K_{1/2}$ by $k$. Its matrix image remains in the normalizer, since conjugation by an element of the torus preserves its normalizer. On complex conjugation the reductions of the two sections are $J,J'$, which are not equal up to sign. This proves distinctness. The attached program verifies every finite identity; at level seven one representative word for $k_7$ is $L_0^{-2}L_1L_0^{-2}$.
\end{proof}

\begin{claim}{D.4}{Proved}{Exact finite restricted splitting counts and a graded obstruction calculation}
The following numbers count sections of the \emph{restriction to $C_2=\langle c\rangle$} of the projective matrix extension $H_n\rtimes\langle c\rangle$, with $H_n$ the image of $F$ generated by $L_0,L_1$ and $c$ acting by conjugation by $J$. They do not count sections over an unspecified ``finite GT group''.
\begin{center}
\begin{tabular}{rrrrrrrrrr}\toprule
$n$&2&3&4&5&6&7&8&12&24\\\midrule
$|H_n|$ projectively&1&12&4&60&12&168&32&96&768\\
Sections over $C_2$&1&6&4&16&6&28&24&48&288\\
Normalizing the $S$-algebra&1&2&4&4&2&4&8&16&32\\\bottomrule
\end{tabular}
\end{center}
For a lower-central graded piece $V_n=\operatorname{gr}_n(F^{(\ell)})\simeq L_n(\mathbb Z_\ell^2)$, complex conjugation acts by $(-1)^n$. Its continuous cohomology is
\[
 H^1(C_2,V_n)=
 \begin{cases}V_n/2V_n& n\text{ odd},\\0&n\text{ even},\end{cases}
 \qquad
 H^2(C_2,V_n)=
 \begin{cases}0&n\text{ odd},\\V_n/2V_n&n\text{ even}.
 \end{cases}
\]
Thus both vanish for odd $\ell$ and over every $\mathbb Q_\ell$. At $\ell=2$ the first possible conjugacy ambiguity is in degree one, of dimension two over $\mathbb F_2$; the first possible lifting obstruction is in degree two, of dimension one. These are cohomological obstruction groups, not assertions that every class is realized by the full extension.
\end{claim}
% anchor:D.4
\manuscriptanchor{\S 46; fols403--403 verso; scans208--209}

\begin{proof}
A section over $C_2$ sends its generator to $(h,c)$, so its necessary and sufficient condition is $h\,c(h)=1$ in $H_n$. Equivalently $(hJ)^2$ is scalar $\pm I$. The program enumerates the finite subgroup generated by $L_0,L_1$, identifies only global $\pm I$, and tests this equality for each element. It then tests whether conjugation sends $S$ into the two-dimensional algebra spanned by $I,S$. Since that conjugation is invertible and the image of $S$ still generates the same algebra in the tested cases, this is the stated normalizer condition. Each displayed number is printed in \texttt{dictionary\_matrices.out}; no inference from these restricted counts to full profinite sections is made.

The two-periodic resolution of $C_2$ has alternating maps $c-1$ and $1+c$. Hence $H^1=\ker(1+c)/\operatorname{im}(c-1)$ and $H^2=\ker(c-1)/\operatorname{im}(1+c)$. Substitute the scalar $(-1)^n$ and use that $V_n$ is torsion-free. The action on the degree-one generators is minus identity; functoriality of a bracket of length $n$ gives the asserted scalar in degree $n$. The Lyndon-word basis proves that these modules are free with Witt ranks, which are computed in the attached script.
\end{proof}

\begin{claim}{D.5}{Proved}{The plus sign is correct in the manuscript's matrix coordinates}
With $R=\left(\begin{smallmatrix}0&1\\-1&1\end{smallmatrix}\right)$, the centralizer in $\operatorname{GL}_2(A)$, for every commutative ring $A$, is
\[
 T_R(A)=\left\{cR+dI:\ c,d\in A,\quad c^2+cd+d^2\in A^\times\right\}.
\]
Thus the plus sign in scan209 (109) agrees with the explicit matrix convention of scans94,154,199 and258. The minus sign is the norm in the basis $(1,\zeta_3)$; the plus sign is the norm in $(1,\zeta_6)$. The manuscript uses $R$ corresponding to $\zeta_6$, not to $\zeta_3$. The parenthetical in the supplied second-pass corrections suggesting that (147)--(150) require a minus sign is mathematically incorrect with this fixed $R$.
\end{claim}
% anchor:D.5
\manuscriptanchor{\S 42; fol336; scan94; \S 46; fol403 verso; scan209(109); \S 48; scans258--259}

\begin{proof}
Writing $M=\left(\begin{smallmatrix}a&b\\c&d\end{smallmatrix}\right)$ and solving $MR=RM$ entry by entry gives $c=-b$, $d=a+b$. Hence $M=aI+bR$, and
$\det(aI+bR)=a^2+ab+b^2$.
There was no division, so the formula holds also in characteristics two and three and over nonreduced rings. The characteristic polynomial is $T^2-T+1$, so its roots are $\zeta_6,\zeta_6^{-1}$ and the displayed determinant is their norm. For $\zeta_3+\zeta_3^{-1}=-1$, the norm of $d+c\zeta_3$ is $d^2-cd+c^2$. Replacing $R$ by $-R$ exchanges the conventions, but that is a change of matrix, not an alternative sign for the same matrix.

The program enumerates all $n^4$ matrices for $n=2,3,4$ and verifies that their invertible centralizer is exactly this norm-unit set; its sizes are $3,6,12$. At $n=3$, $c=d=1$ gives a singular matrix: the correct norm is $0$, whereas the minus-sign expression is $1$. At $n=2$ the two expressions coincide, so level two alone cannot decide the reading. At $n=4$ their values can differ although their unit/nonunit parity is the same. At the prime three this smooth centralizer group is not a torus: the algebra $\mathbb F_3[R]$ is the dual-number algebra because $T^2-T+1=(T+1)^2$; its units have a nontrivial additive subgroup. This is exactly the caveat Grothendieck records in \S48, scan259.
\end{proof}

\begin{claim}{D.6}{Open}{The remaining order-three and global splitting problem}
The work above does not classify sections of $\widetilde M^!$ over all of $\widehat{\mathrm{GT}}_0$ up to $\widehat F$-conjugacy, nor does it classify sections over $C(3)=\ker(\mu\bmod3)$ whose matrix image lies in $T_R$. The exact remaining question is whether the subset of
$H^1_{\mathrm{cts}}(C(3),\widehat F)$ satisfying that torus constraint, for the action of the centralizing section, consists of one class. For a compatible lower-central or finite characteristic quotient system, a lift at the next abelian kernel $V$ is obstructed by its explicit extension class in $H^2_{\mathrm{cts}}(C(3),V)$; if zero, choices form an $H^1$-torsor. Already $V=\widehat F^{\mathrm{ab}}=\widehat{\mathbb Z}(1)^2$ is an undecided global instance in this return. D.4 computes the full analogous obstruction groups after restriction to the involution subgroup; that subgroup is not contained in $C(3)$ and is not a substitute for it.
\end{claim}
% anchor:D.6
\manuscriptanchor{\S 46; fol403 verso; scan209}

\begin{proof}[Established reduction and attempted attack]
D.1 proves existence for the whole extension without a torus constraint, and the $C_3$-centralizer construction restricts its domain as in B.m4. For any successive extension $1\to V\to E'\to E\to1$, pull it back along the proposed section of $E$. Choose a continuous local set-theoretic lift $b$ when one exists and form
$\kappa(g,h)=b(g)b(h)b(gh)^{-1}$.
Associativity gives the two-cocycle equation. Changing the lift changes $\kappa$ by a coboundary, so a homomorphic lift exists precisely when its class vanishes; differences of lifts satisfy the one-cocycle equation. This is the stated obstruction, and the finite restricted calculation tests its first involution instances explicitly. The failed further step is to compute this class and its constrained $H^1$ for the full profinite domain $C(3)$, not to multiply a longer list of $2\times2$ matrices. A finite quotient of the target alone supplies neither that domain nor compatible transition sections. We leave the exact global instance visible rather than infer uniqueness from the successful rational averaging argument.
\end{proof}

\begin{claim}{D.7}{Refuted}{Strict uniqueness with the order-three torus constraint also fails}
Let $C(3)=\ker(\mu\bmod3)$ and let $s_R:C(3)\to\widetilde M^!$ be the section centralizing $r$. There is a distinct conjugate section whose matrix image still lies in $T_R$; indeed the two sections have identical images at every congruence level. Their difference is detected in the characteristic abelian quotient $\widehat F^{\rm ab}/9$. Thus the strict uniqueness suggested and withdrawn on \S46, fol.403$'$, scan209 is false even with its stated order-three torus constraint. This does not decide uniqueness modulo geometric conjugation.
\end{claim}
\begin{proof}
For $\mu\equiv1\pmod3$, the conjugator of $u(r)$ may be chosen in $\widehat F$, by the factor-torsion argument in C.d1. Its ambiguity is $C_{\widehat B}(r)\cap\widehat F=\langle r\rangle\cap\widehat F=1$. Adjusting by that unique conjugator gives a unique representative centralizing $r$; these representatives form the continuous section $s_R$.

Take the determinant-one adelic matrix $k^{\rm ad}$ with component $R^2$ at three and $I$ at every other prime. It reduces to $I$ modulo two and centralizes $R$, so the geometric surjectivity proved in D.3 gives $k\in\widehat F$ with that matrix image. In $\operatorname{SL}_2(\mathbb F_3)/\{\pm I\}\simeq A_4$, the image of $R^2$ has order three and is nontrivial in $A_4^{\rm ab}\simeq C_3$. Therefore $[k]\bmod3\ne0$ in $\widehat F^{\rm ab}/3$. The program independently computes this derived subgroup and tests membership.

Choose the unit $\mu\in\widehat{\mathbb Z}^{\times}$ with $\mu_3=4$ and $\mu_p=1$ for $p\ne3$. Surjectivity of the cyclotomic character gives an arithmetic element and hence $g\in C(3)$ with this multiplier. Since every pure loop automorphism acts on $\widehat F^{\rm ab}$ by its multiplier, the cocycle comparing $s_R$ and $k s_R k^{-1}$ at $g$ has abelianization
\[
 [k]-\mu[k]=(1-\mu)[k],
\]
whose three-adic component is $-3[k]$ and is nonzero modulo nine. Thus the two sections are strictly distinct. The finite calculation checks this for every vector modulo nine whose reduction modulo three is nonzero.

On the other hand, $U(s_R(h))$ centralizes $R$ for every $h\in C(3)$. A possible global negative sign in the projective equality is ruled out by $\operatorname{tr}R=1$. The centralizer algebra $\widehat{\mathbb Z}[R]$ is commutative, so $k^{\rm ad}$ commutes with all these matrices. Consequently conjugating the section by $k$ changes none of its congruence matrices. Both sections satisfy the exact torus constraint. This explains why a search solely in the congruence quotients would miss the counterexample, although the abelian characteristic quotient detects it.
\end{proof}

